\cotefile{136}
\scanpage{267}
Si $\deleted{\unsure{i}}\unsure{\mathcal C|C}\xrightarrow{i}\mathcal M_{N_0}$ est un foncteur exact

\[
\operatorname{Im}\deleted{\unsure{\mathrm{ess}}}\,i\subset\operatorname{Im\,ess}\varphi_{N_0,\infty}
\]

alors on a $i\simeq\varphi_{N_0,\infty}\tau_{\infty,N_0}i$ d'où

\[
(\mathbb R^{\deleted{\unsure{*|n}}}\varphi_{N_0,\infty})(\tau_{\infty,N_0}\circ i)\simeq\mathbb R^*i
\]

en particulier

\[
(R^n\varphi_{N_0,\infty})(\tau_{\infty\,N_0}i(M))=0
\qquad\begin{gathered}\text{si }n>0\\ M\in\mathrm{ob}\,C.\end{gathered}
\]

On est intéressé au = cas particulier qui avant.

\figure{\begin{tikzpicture}\draw (0,0)--(1.45,0);\end{tikzpicture}}

\[
\begin{aligned}
\text{Soit}\qquad &K^\cdot\text{ un complexe de }C\\
&L^\cdot\text{ un complexe de }\mathcal M_{\deleted{\unsure{\infty}}}
\end{aligned}
\]

et supposons donné un hom dans la catégorie dérivée $D(\mathcal M_{N_0})$

\[
i(K^\cdot)\xrightarrow{\alpha}L^\cdot
\]

tel que

\[
\tau_{\infty,N_0}i(K)\xrightarrow[\sim]{\tau_\infty\alpha}\tau_{\infty,N_0}L^\cdot
\]

soit un isom dans $D(\mathcal M_\infty)$, i.e. tel que $\forall n$,

\[
H^n(\alpha):\underset{\substack{\wr\\ i(H^n(K))}}{H^{n}\deleted{\unsure{\varphi}}(i(K))}\longrightarrow H^n(L^\cdot)
\]

soit un $\infty$-isom, i.e. noyau et conoyau \unsure{sont|soit} négligeables. On en conclut

% archive sheet number: 266
